\documentclass[11pt,a4paper]{article}

% --- Encoding and fonts ---
\usepackage[utf8]{inputenc}
\usepackage[T1]{fontenc}

% --- Geometry: wider text block + tighter top margin (no title page) ---
\usepackage[top=1.8cm,bottom=2cm,left=2cm,right=2cm]{geometry}

% --- URL line-breaking BEFORE hyperref so DOIs in citations wrap ---
\usepackage[hyphens]{url}

% --- Bibliography (biblatex + biber) ---
\usepackage[backend=biber,style=authoryear,sorting=nyt,maxbibnames=6]{biblatex}
% \runDir is injected via latexmk -usepretex='\def\runDir{runs/<run_id>}'
\IfFileExists{\runDir/references.bib}{%
  \addbibresource{\runDir/references.bib}%
}{}

% --- Tables, figures, links ---
\usepackage{longtable}
\usepackage{booktabs}
\usepackage{array}
\usepackage{graphicx}
\usepackage[colorlinks=true,linkcolor=blue,citecolor=blue,urlcolor=blue,breaklinks=true]{hyperref}

% --- Math ---
\usepackage{amsmath}

% --- Robust line breaking for long unbreakable strings ---
\setlength{\emergencystretch}{3em}
\hbadness=10000
\sloppy

% --- Default longtable to footnotesize so wide tables stay in margins ---
\usepackage{etoolbox}
\AtBeginEnvironment{longtable}{\footnotesize}

% --- Cap any \includegraphics that pandoc emits at \textwidth. ---
\setkeys{Gin}{width=\textwidth,keepaspectratio}

% --- Pandoc 3.x wraps figures with \pandocbounded{...}; provide a no-op stub. ---
\providecommand{\pandocbounded}[1]{#1}

% --- Pandoc emits \tightlist for compact bullet lists; provide a stub. ---
\providecommand{\tightlist}{\setlength{\itemsep}{0pt}\setlength{\parskip}{0pt}}

% --- Auto-cite all entries in references.bib so the bibliography is
%     populated even though we don't \cite{} keys inline. ---
\nocite{*}

% --- PDF metadata: avoid leaking the local run path in document properties. ---
\hypersetup{
  pdftitle={LitChron Annotation Report},
  pdfauthor={LitChron MCP Server},
  pdfsubject={Single-cell pseudotime annotation},
  pdfkeywords={LitChron, single-cell, pseudotime, LLM, annotation},
}

\begin{document}

% Compact heading instead of a dedicated title page — no whitespace waste,
% no filesystem path exposure.
\begin{center}
{\Large\bfseries LitChron Annotation Report}\\[2pt]
{\small\itshape Single-cell pseudotime annotated by LLM, supporting
literature verified against CrossRef}\\[1pt]
{\footnotesize \today}
\end{center}
\vspace{4pt}

% 1. Observations — cluster markers from the recomputed embedding.
\IfFileExists{\runDir/tex_sections/observations.tex}{%
  \input{\runDir/tex_sections/observations}%
}{%
  \section{Dataset Observations}\textit{Observations section not yet generated.}%
}

% 2. LitChron Annotation — the headline single-panel UMAP figure.
\IfFileExists{\runDir/tex_sections/annotation.tex}{%
  \input{\runDir/tex_sections/annotation}%
}{}

% 3. References — placed immediately after the figure so the supporting
%    literature sits with the biological claims it underpins.
\printbibliography[heading=bibintoc,title={References}]

% 4. LLM Proposal — per-cell-type ranking table.
\IfFileExists{\runDir/tex_sections/proposals.tex}{%
  \input{\runDir/tex_sections/proposals}%
}{%
  \section{LLM-Proposed Pseudotime Ordering}\textit{Proposals section not yet generated.}%
}

% 5. LitChron LLM Pseudotime — per-cell continuous pseudotime summary.
\IfFileExists{\runDir/tex_sections/litchron_ordering.tex}{%
  \input{\runDir/tex_sections/litchron_ordering}%
}{}

\end{document}
